\begin{pastexam}[label={px:ncku-114-ee-ds-17}]{成大電機所 114 學年度 簡答題第 1 題}
（填空）

(10 PTs) Given a circular list (or, a singly-linked circular list) that the link field of the last
node points to the first node as shown in Fig 1, please complete the code to insert a new node to
the front of the list in the following:

\begin{minipage}[t]{0.55\textwidth}
\begin{minted}[fontsize=\footnotesize,escapeinside=||]{cpp}
template <class T>
void CircularList<T>::InsertFront(const T& e)
{
  ChainNode <T> *n = new ChainNode <T>(e);
  if (last) { // nonempty chain
     |\blank{a)}|;
     |\blank{b)}|;
  } else {
     |\blank{c)}|;
     |\blank{d)}|;
  }
}
\end{minted}
\end{minipage}\hfill
\begin{minipage}[t]{0.4\textwidth}
\begin{minted}[fontsize=\footnotesize]{cpp}
template <class T> class Chain;

template <class T>
class ChainNode {
  friend class Chain <T>;
  private:
    T data;
    ChainNode <T> *link;
};
\end{minted}
\end{minipage}

\begin{center}
\begin{tikzpicture}[>={Stealth[length=2.2mm]},font=\small]
  \foreach \k/\x in {1/0,2/3.2,3/6.4} {
    \draw[fill=black!25] (\x,0) rectangle (\x+1.7,0.6);
    \draw (\x+1.7,0) rectangle (\x+2.4,0.6);
    \node[inner sep=1pt] at (\x+0.85,0.3) {$x_{\k}$};
  }
  \draw[->] (2.05,0.3) -- (3.2,0.3);
  \draw[->] (5.25,0.3) -- (6.4,0.3);
  \draw[->] (8.45,0.3) -- (9.2,0.3) -- (9.2,1.2) -- (0.85,1.2) -- (0.85,0.6);
  \node[font=\small\sffamily\itshape] at (0.85,-0.3) {data};
  \node[font=\small\sffamily\itshape] at (2.05,-0.3) {link};
  \draw[->] (8.45,-0.6) -- (8.45,-0.05);
  \node[font=\small\sffamily\itshape] at (8.45,-0.85) {last};
  \node at (4.4,-0.85) {Fig. 1};
\end{tikzpicture}
\end{center}
\end{pastexam}

\begin{pastsoln}
答案：
\begin{minted}[fontsize=\small]{cpp}
if (last) {                // 非空：接在 last 後面，也就是第一個節點前面
    n->link = last->link;  // a) 新節點指向原本的第一個節點
    last->link = n;        // b) 最後一個節點改指向新節點
} else {                   // 空串列：新節點自己成一個環
    last = n;              // c)
    n->link = n;           // d)
}
\end{minted}

環狀串列記的是最後一個節點 \texttt{last}，第一個節點是 \texttt{last->link}。 在最前面插入，就是把
新節點接在 \texttt{last} 和原本的第一個節點之間；\texttt{last} 不用動，因為最後一個節點沒變
（資結書第~5.2~節，和虛擬碼~5.5 在最後面插入只差最後要不要移 \texttt{last}）。

a)、b) 的順序不能對調：先做 \texttt{last->link = n} 的話，原本的第一個節點就找不到了。 空串列時新
節點既是第一個也是最後一個，\texttt{link} 指向自己；c)、d) 兩行的順序可以對調（寫成
\texttt{last = n; last->link = last;} 也可以）。 例如依序在空串列前端插入 $1, 2, 3$，得到環
$3 \to 2 \to 1 \to 3$，\texttt{last} 一直指向 $1$。 時間 $\Theta(1)$。

附帶一提：右邊的 \texttt{ChainNode} 只把 \texttt{Chain<T>} 宣告成 friend，要讓
\texttt{CircularList<T>} 存取 private 的 \texttt{link}，實際程式還得把 \texttt{CircularList<T>}
也宣告成 friend。 這不影響填空的答案。
\end{pastsoln}
